% ------------------------------------------------------------------
% HW1 Problem 2(f) -- derivation of the cosine-wave velocity commands
% Paste this into the Task 2 section of the solution file.
% Needs: \usepackage{amsmath}
% ------------------------------------------------------------------
\subsection*{Problem 2(f): Driving the turtle on a cosine wave}

\paragraph{Model.} The turtle has differential-drive (unicycle) kinematics. It can only be
commanded a forward speed $v$ along its body $x$-axis and a turn rate $\omega$ about $z$,
i.e.\ \texttt{linear.x} and \texttt{angular.z} of the \texttt{Twist} on
\texttt{/turtle1/cmd\_vel}. In the world frame,
\begin{equation}
  \dot{x} = v\cos\theta, \qquad \dot{y} = v\sin\theta, \qquad \dot{\theta} = \omega .
  \label{eq:unicycle}
\end{equation}
The task is the inverse problem: choose $v(t)$ and $\omega(t)$ so that the turtle follows a
prescribed path.

\paragraph{Desired path.} Move to the right at a constant rate $c$ while oscillating in $y$:
\begin{equation}
  x(t) = c\,t, \qquad y(t) = A\cos(kct),
\end{equation}
so the trace is $y = A\cos(kx)$. At $t=0$ the slope is zero, which matches the turtle's
default heading $\theta = 0$.

\paragraph{Derivatives.}
\begin{align}
  \dot{x} &= c, & \ddot{x} &= 0, \\
  \dot{y} &= -Akc\,\sin(kct), & \ddot{y} &= -Ak^{2}c^{2}\cos(kct).
\end{align}

\paragraph{Forward speed.} Squaring and adding the first two equations of \eqref{eq:unicycle}
gives $\dot{x}^2+\dot{y}^2 = v^2$, so
\begin{equation}
  v(t) = \sqrt{\dot{x}^{2}+\dot{y}^{2}}
       = \sqrt{c^{2} + A^{2}k^{2}c^{2}\sin^{2}(kct)}
       = c\sqrt{1 + A^{2}k^{2}\sin^{2}(kct)} .
\end{equation}

\paragraph{Heading and turn rate.} Dividing the same two equations gives
$\tan\theta = \dot{y}/\dot{x}$, so $\theta = \arctan(u)$ with
\begin{equation}
  u = \frac{\dot{y}}{\dot{x}} = -Ak\sin(kct), \qquad
  \dot{u} = \frac{\ddot{y}}{c} = -Ak^{2}c\cos(kct), \qquad
  u^{2} = A^{2}k^{2}\sin^{2}(kct).
\end{equation}
Using $\frac{d}{dt}\arctan(u) = \dfrac{\dot{u}}{1+u^{2}}$,
\begin{equation}
  \omega(t) = \dot{\theta} = \frac{\dot{u}}{1+u^{2}}
  = \frac{-Ak^{2}c\,\cos(kct)}{1 + A^{2}k^{2}\sin^{2}(kct)} .
\end{equation}

\paragraph{Checks.}
\begin{itemize}
  \item Top of the wave, $t=0$: $\sin = 0$, $\cos = 1$, so $v = c$ and
        $\omega = -Ak^{2}c < 0$, the hardest right turn, since the path bends down.
        The turning radius there is $R = v/\omega = -1/(Ak^{2})$, i.e.\ the ICR is
        $1/(Ak^{2})$ to the turtle's right.
  \item Steepest point, $kct = \pi/2$: $\cos = 0$, so $\omega = 0$ (the bend changes
        direction) and $v = c\sqrt{1+A^{2}k^{2}}$ is largest, because the turtle covers
        the most ground there.
\end{itemize}

\paragraph{Implementation and result.} Node \texttt{turtle\_wave/cosine\_wave} with
$A=2$, $k=1$, $c=1$, turtle teleported to $(0.5,\,7.0,\,0)$.
\begin{enumerate}
  \item \emph{Open loop} ($v(t)$, $\omega(t)$ published at 50\,Hz from a clock): the first
        half-wave was correct, then the path drifted and curled. Nothing corrects late or
        missed commands, so heading errors accumulate, as with encoder-only dead reckoning.
  \item \emph{With pose feedback}: the same $v$ and $\omega$ are used as feed-forward, evaluated
        at the turtle's measured position ($kct \to k(x-x_0)$), plus a correction
        $\omega = \omega_{\mathrm{ff}} - K_\theta(\theta-\theta_{\mathrm{des}}) - K_y(y-y_{\mathrm{des}})$,
        $K_\theta = 4$, $K_y = 2$, using \texttt{/turtle1/pose}. The turtle traced a clean
        cosine of about 1.5 periods between $y=3$ and $y=7$.
\end{enumerate}

\noindent\textit{AI use: Claude used for the ROS node boilerplate, build steps, the
pose-feedback correction, and typesetting this derivation in \LaTeX.}
